\section{Macroeconomía: contabilidad nacional, inflación, dinero, desempleo y política fiscal}
\label{sec:macroeconomia}

La macroeconomía estudia el funcionamiento de la economía en su conjunto. A diferencia de la microeconomía, que analiza individuos y mercados específicos, la macroeconomía observa agregados, como el ingreso total, la inflación, el desempleo, la producción nacional y la política económica del Estado.

\subsection{Contabilidad nacional}

La contabilidad nacional mide la actividad económica de un país a nivel agregado. La variable más importante es el Producto Interno Bruto (PIB).

\begin{concepto}{PIB}
El PIB es el valor de mercado de todos los bienes y servicios finales producidos en un país durante un período determinado.
\end{concepto}

\begin{equation}
Y = C + I + G + XN
\end{equation}

donde:
\begin{itemize}
    \item $Y$: PIB
    \item $C$: consumo
    \item $I$: inversión
    \item $G$: gasto público
    \item $XN$: exportaciones netas
\end{itemize}

\begin{ejemplo}
Si un país tiene consumo por $120$, inversión por $50$, gasto público por $40$ y exportaciones netas por $10$, entonces su PIB es:
\[Y = 120 + 50 + 40 + 10 = 220
\]
\end{ejemplo}

\subsection{PIB nominal y PIB real}

\begin{itemize}
    \item \textbf{PIB nominal:} mide la producción a precios corrientes del período.
    \item \textbf{PIB real:} ajusta el valor por inflación para ver cuánto cambió la producción en términos reales.
\end{itemize}

\begin{observacion}{Importante}
El PIB nominal puede aumentar por dos razones: porque se produce más o porque los precios suben. El PIB real separa estos dos efectos.
\end{observacion}

\subsection{PIB per cápita}

\begin{concepto}{PIB per cápita}
Es el PIB dividido por la población del país.
\end{concepto}

Es una medida útil para comparar niveles de ingreso promedio entre países, aunque no captura completamente bienestar o distribución.

\begin{ejemplo}
Si un país tiene PIB de $500$ mil millones y una población de $50$ millones, entonces:
\[\text{PIB per cápita} = \frac{500\,000\,000.000}{50\,000.000} = 10.000
\]
\end{ejemplo}

\subsection{Flujo circular de la renta}

La economía funciona como un flujo circular entre hogares y empresas:
\begin{itemize}
    \item las familias venden trabajo y otros recursos a las empresas,
    \item las empresas pagan salarios y rentas,
    \item las familias gastan ese ingreso en bienes y servicios,
    \item las empresas producen esos bienes y servicios.
\end{itemize}

\begin{ejemplo}
Un trabajador recibe sueldo, lo usa para comprar comida y transporte, y esas compras permiten a las empresas seguir produciendo y contratar más trabajo.
\end{ejemplo}

\subsection{Productividad y crecimiento económico}

El crecimiento económico depende de varias variables, como:
\begin{itemize}
    \item innovación tecnológica,
    \item capital humano,
    \item inversión,
    \item infraestructura,
    \item instituciones estables,
    \item educación y salud.
\end{itemize}

\begin{concepto}{Productividad}
Es la cantidad de producción que se obtiene por unidad de trabajo o recursos.
\end{concepto}

\begin{ejemplo}
Si una fábrica mejora su maquinaria y produce más con la misma cantidad de trabajadores, su productividad aumenta.
\end{ejemplo}

\subsection{IPC e inflación}

\begin{concepto}{Inflación}
Es el aumento sostenido del nivel general de precios de la economía.
\end{concepto}

El Índice de Precios al Consumidor (IPC) mide el costo de una canasta de bienes y servicios representativa del consumo de una familia promedio.

\begin{equation}
\text{IPC} = \frac{\text{costo de la canasta actual}}{\text{costo de la canasta base}} \times 100
\end{equation}

\begin{ejemplo}
Si la canasta base costaba $100$ y ahora cuesta $110$, el IPC es $110$ y la inflación entre ambos períodos es $10\%$.
\end{ejemplo}

\subsection{Inflación: causas y efectos}

La inflación puede deberse a diferentes causas:
\begin{itemize}
    \item aumento de la demanda agregada,
    \item aumento de costos de producción,
    \item crecimiento excesivo del dinero,
    \item choques de oferta.
\end{itemize}

\begin{ejemplo}
Si el gobierno imprime demasiado dinero, la cantidad de dinero en circulación aumenta y los precios se elevan. Esto es inflación monetaria.
\end{ejemplo}

\begin{advertencia}{Importante}
No toda subida de precios en un bien aislado significa inflación general. La inflación se refiere al aumento general del nivel de precios de la economía.
\end{advertencia}

\subsection{Deflactor del PIB frente al IPC}

\begin{itemize}
    \item \textbf{IPC:} mide precios de bienes de consumo.
    \item \textbf{Deflactor del PIB:} mide el cambio de precios del conjunto de bienes producidos en la economía.
\end{itemize}

\begin{ejemplo}
Si el IPC sube porque aumentan los precios del transporte y la alimentación, eso refleja inflación del consumo; si el deflactor del PIB también sube, hay inflación más general en producción.
\end{ejemplo}

\subsection{Dinero: funciones y tipos}

El dinero tiene varias funciones:
\begin{enumerate}
    \item medio de intercambio,
    \item unidad de cuenta,
    \item depósito de valor,
    \item reserva de valor.
\end{enumerate}

\begin{concepto}{Dinero}
Es cualquier bien aceptado como medio de pago en una economía.
\end{concepto}

\begin{ejemplo}
El dinero en efectivo, depósitos bancarios y transferencias electrónicas cumplen funciones monetarias en la economía moderna.
\end{ejemplo}

\subsection{Banco Central}

El Banco Central es la autoridad monetaria del país. Su función principal es mantener la estabilidad monetaria y financiera.

\begin{concepto}{Banco Central}
Es la institución encargada de la política monetaria y del control del sistema financiero.
\end{concepto}

Sus funciones incluyen:
\begin{itemize}
    \item regular el sistema bancario,
    \item mantener la inflación bajo control,
    \item influir en las tasas de interés,
    \item ofrecer liquidez al sistema financiero.
\end{itemize}

\subsection{Política monetaria}

\begin{concepto}{Política monetaria}
Es el conjunto de medidas del Banco Central para influir en la cantidad de dinero, el crédito y las tasas de interés.
\end{concepto}

\begin{itemize}
    \item Política monetaria expansiva: baja tasas de interés para estimular la actividad.
    \item Política monetaria contractiva: sube tasas de interés para frenar inflación o sobrecalentamiento.
\end{itemize}

\begin{ejemplo}
Si el Banco Central sube la tasa de interés, el crédito se encarece y la actividad económica tiende a desacelerarse.
\end{ejemplo}

\subsection{Sistema bancario y creación de dinero}

Los bancos no solo guardan depósitos: también crean crédito y, por lo tanto, dinero bancario. El sistema bancario puede ampliar la oferta monetaria mediante préstamos.

\begin{observacion}{Importante}
No todo el dinero se crea físicamente; una gran parte del dinero en circulación es dinero bancario, generado mediante préstamos.
\end{observacion}

\subsection{Desempleo}

\begin{concepto}{Desempleo}
Es la situación de quienes quieren trabajar y están disponibles para hacerlo, pero no encuentran empleo.
\end{concepto}

Tipos de desempleo:
\begin{itemize}
    \item \textbf{friccional:} entre trabajos o búsqueda inicial,
    \item \textbf{estructural:} por desajuste entre habilidades y vacantes,
    \item \textbf{cíclico:} asociado a recesiones.
\end{itemize}

\begin{ejemplo}
Una persona que cambia de trabajo puede pasar un tiempo desempleada por elección o búsqueda; eso es desempleo friccional. Si hay crisis económica y muchos despidos, el desempleo cíclico aumenta.
\end{ejemplo}

\subsection{Política fiscal}

La política fiscal es el conjunto de decisiones del Estado sobre gasto público e impuestos.

\begin{itemize}
    \item \textbf{Política fiscal expansiva:} aumenta gasto o baja impuestos.
    \item \textbf{Política fiscal contractiva:} reduce gasto o sube impuestos.
\end{itemize}

\begin{ejemplo}
Durante una crisis, el Estado puede aumentar gasto público en infraestructura para sostener la actividad económica.
\end{ejemplo}

\subsection{Deuda pública}

\begin{concepto}{Deuda pública}
Es el conjunto de obligaciones financieras que tiene el Estado con terceros.
\end{concepto}

La deuda puede ser útil si se usa para invertir en infraestructura, educación o salud, pero puede volverse problemática si crece sin capacidad de pago.

\begin{advertencia}{Importante}
No toda deuda es mala ni toda deuda es buena. Lo importante es su nivel, su uso y la capacidad del país para sostenerla.
\end{advertencia}

\subsection{Conclusión macroeconómica}

La macroeconomía busca entender cómo funciona la economía en conjunto: cómo se mide la producción, cómo cambia el nivel de precios, cómo se crea dinero, cómo se mide el desempleo y cómo reaccionan los gobiernos con política fiscal y monetaria.

\begin{resumen}{Puntos clave de la macroeconomía}
- El PIB mide la producción agregada.
- La inflación mide el aumento general de precios.
- El dinero cumple funciones esenciales para la economía.
- El Banco Central controla la política monetaria.
- El desempleo es un indicador central del bienestar económico.
- La política fiscal y la deuda pública determinan parte del equilibrio macroeconómico.
\end{resumen}
